\subsection{Basic Facts About Modular Arithmetic}

\content{
        \begin{example} Compute the following:
        \begin{enumerate}
        \item $ 47^{8} \mod{49} $
        \item $ 65^{5} \mod{35} $
        \end{enumerate}
        \end{example}
}{% presentation
\vspace{3in}
}{% nopresentation
\vspace{3in}
}{% comments
}

\content{
        \begin{example} Your turn! Compute the following:
        \begin{enumerate}
        \item $ 12^8 \mod{4} $
        \item $ 12^8 \mod{5} $
        \item $ 5^{12} \mod{4} $
        \item $ 98^{123} \mod{99} $
        \end{enumerate}
        \end{example}
}{% presentation
}{% nopresentation
}{% comments
}

\content{
        What else can we do with modular arithmetic? 

        \begin{example} Compute each of the following:
        \begin{enumerate}
        \item $3^{10} \mod{22}$
        \item $5^{10} \mod{22}$
        \item $21^{10}\mod{22}$
        \end{enumerate}
        \end{example}
}{% presentation
\vspace{3in}
}{% nopresentation
\vspace{3in}
}{% comments
Talk briefly about Euler's Theorem: 
$$a^{\phi(n)}\equiv 1\mod{n}$$
when $\gcd(a,n)=1$.
}

\content{
        \begin{definition} Euler's totient function $\phi$ is defined as
        follows: If $n$ is a positive whole number, and $p_1,p_2,\dots,p_n$ are
        all the prime numbers dividing $n$, then 
        $$ \phi(n) =
        n\left(1-\frac{1}{p_1}\right)\left(1-\frac{1}{p_2}\right)\cdots
        \left(1-\frac{1}{p_n}\right)$$
        \end{definition}
}{% presentation
}{% nopresentation
}{% comments
}

\content{
        \begin{example} Compute the following:
        \begin{enumerate}
        \item $\phi(27)$
        \item $\phi(19)$
        \item $\phi(42)$
        \item $\phi(124)$
        \end{enumerate}
        \end{example}
}{% presentation
\vspace{2in}
}{% nopresentation
\vspace{2in}
}{% comments
}

\content{
        \begin{example} So how might we compute something like the following:
        $$ 128^{1324} \mod{34} $$
        \end{example}
}{% presentation
\vspace{2in}
}{% nopresentation
\vspace{1in}
}{% comments
16
}

\content{
        \begin{example} Your turn! Compute
        $$ 415^{1521} \mod{27}. $$
        \end{example}
}{% presentation
\vspace{1in}
}{% nopresentation
\vspace{1in}
}{% comments
1
}


\subsubsection{Key Exchange}

\content{
        Up to this point every method of sending an encrypted message requires a
        secret key known to both Alice and Bob. A big problem in cryptography is
        the following: 

        Suppose that Alice and Bob are unable to meet in person, they have to
        email each other to agree on the secret key.  Unfortunatly, Eve has
        access to Alice's emails, and thus can read any emails that are sent
        between Alice and Bob. How can Alice and Bob decide on a secret key
        without the key falling into Eve's hands? 
}{% presentation
\newpage
}{% nopresentation
\vspace{3in}
}{% comments
We need some `one-way function' that is easy to compute, but hard to undo. 

Illustrate how this can be done with the paint mixing example. In this example,
the secret key is the final paint color. 

Of course, since Alice and Bob are using computers, the one-way function will
need to be a numerical function. 
}
